\cotefile{63}
\scanpage{3}
\begin{center}\uline{Contents}\end{center}

\noindent 1. Generalized Weierstrass form \qquad p~1.

\noindent 2. Change of coördinates \qquad p~2.

\noindent 3. \unsure{'}Minimal \unsure{weierstrass|Weierstrass} equations over a valuation ring \qquad p~3.

\noindent 4. The canonical filtration on the group of $v$-adic points \qquad p~4.

\noindent 5. Application\unsure{.|:} The relation between $L_v(1)$ and $\int_{E_v}|\omega_v|$ in case $k$ is finite \qquad p~6.

\noindent 6. The Neron minimum model \qquad p~7.

\noindent 7\unsure{,|.} Algorithm for analysing singular fibers \quad (first five cases) \qquad p~8.

\noindent 8. Algorithm continued \quad (last five cases) \qquad p~10.

% archive sheet number: 2
\scanpage{4}\mspage{1}
\begingroup
\makeatletter\tagsleft@true\makeatother
1. \uline{Generalized} \uline{Weierstrass} \uline{form}. Let $E$ be an \unsure{elliptic|elleptic} curve defined over a field \unsure{-} $K$ with a $K$-rational point $O$. In the projective embedding defined by $3O$ the curve can be written in the form
\[
y^2+a_1xy+a_3y=x^3+a_2x^2+a_4x+a_6
\qquad \unsure{\ldotp|,}\quad\unsure{,}\quad a_{\unsure{i|1}}\in K
\tag*{(1.1)}
\]

Put
\[
\begin{aligned}
 b_2&=a_1^2+4a_2                  & c_4&=b_2^2-24b_4\\
 b_4&=a_1a_3+2a_4                  & c_6&=-b_2^3+36b_2b_4-216b_6\\
 b_6&=a_3^2+4a_6                    & \Delta&=-b_2^2b_8-8b_4^3-27b_6^2+9b_2b_4b_6\ne0\\
 b_8&=a_1^2a_6-a_1a_3a_4+4a_2a_6+a_2a_3^2-a_4^2
                                      & j&=c_4^3/\Delta
\end{aligned}
\tag*{(1.2)}
\]

These quantities are related by
\[
4b_8=b_2b_6-b_4^2 \qquad 1728\Delta=c_4^3-c_6^2
\tag*{(1.3)}
\]

A differential of first kind is given by
\[
\omega=\frac{dx}{2y+a_1x+a_3}
=\frac{dx}{F_y(x,y)}=-\frac{dy}{F_x(x,y)}
=\frac{dy}{3x^2+2a_2x+a_4-a_1y},
\tag*{(1.4)}
\]
where we have put
\[
F(X,Y)\unsure{=|-}Y^2+a_1XY+a_3Y-X^3-a_2X^2-a_4X-a_6.
\tag*{(1.5)}
\]

Putting
\[
\eta=y+\frac{a_1x+a_3}{2}\qquad\text{and}\qquad
\xi=x+\frac{b_2}{12},
\tag*{(1.6)}
\]
the equation (1.1) becomes
\[
\eta^2=x^{\unsure{3|2|1}}+\frac{b_2}{4}x^2+\frac{b_4}{2}x+\frac{b_6}{4}
=\xi^{\unsure{3|2}}-\frac{c_4}{48}\xi-\frac{c_6}{864}.
\tag*{(1.7)}
\]

Thus the relation with Weierstrass is given by
\[
\begin{aligned}
\xi&=\wp(u) & {\unsure{C|c}}_4&=12g_2 &
\omega&=\frac{\unsure{d}\xi}{2\eta}=du\\
2\eta&\unsure{-|=}\wp^{\unsure{\prime}}(u) & C_6&=216g_3
\end{aligned}
\tag*{(1.8)}
\]
\endgroup

% archive sheet number: 3
\scanpage{5}\mspage{2}
\begingroup
\makeatletter\tagsleft@true\makeatother
2. \uline{Change} \uline{of} \uline{coordinates}. Suppose $E':\quad y^{\prime\,2}+a_1'x'y'+\ldotp\,\ldotp\,\ldotp\,\ldotp$ is another curve of the same type as $E$, and $f:E'\overset{\sim}{\longrightarrow}E$ an \unsure{isomorphism} over $K$\unsure{.} carrying $O'$ into $O$. Then there are $r,s,t$ and $u\ne0$ in $K$ such that
\[
\begin{aligned}
x\circ f&=u^2x'+r & \omega\circ f&=u^{-1}\omega'.\\
y\circ f&=u^3y'+su^2x'+t
\end{aligned}
\tag*{(2.1)}
\]

The coefficients $a_i'$ are related to the $a_i$ \deleted{\illegible{1}} and the $b_i'$ to the $b_i$ by the formulas:
\begingroup\scriptsize
\[
\begin{aligned}
ua_1'&=a_1+2s & u^8b_8'&=b_8+3rb_6+3r^2b_4+r^3b_2+3r^{\unsure{4}}\\
u^2a_2'&=a_2-sa_1+3r-s^2 & u^6b_6'&=b_6+2rb_4+r^2b_2+4r^3\\
u^3a_3'&=a_3+ra_1+2t=F_y(r,t) & u^4b_4'&=b_4+rb_2+6r^2\\
u^4a_4'&=a_4-sa_3+2ra_2-(t+rs)a_1+3r^2-2st
 & u^2b_2'&=b_2+12r.\\
 &= -F_x(r,t)-sF_y(r,t)\\
u^6a_6'&=a_6+ra_4+r^2a_2+r^3-ta_3-t^2-rta_1\\
 &= -F(r,t)
\end{aligned}
\tag*{(2.2)}
\]
\endgroup

From these we check that:
\[
u^4c_4'=c_4\qquad u^6c_6'=c_6\qquad
u^{12}\Delta'=\Delta\qquad j'=j.
\tag*{(2.3)}
\]

\uline{Example} \unsure{`|``}generic $E_j$''. The equation
\[
y^2+xy=x^3-\frac{36}{j-1728}x-\frac{1}{j-1728}
\tag*{(2.4)}
\]
has $c_4=c_6=\dfrac{j}{j-1728}$ and $\Delta=\dfrac{j^2}{(j-1728)^3}$. Hence for $j\ne0,1728$ it gives a curve with invariant \unsure{'}$j$.

Using the formulas above, it is easy to show that $j$ can be arbitrary in $K$, and that for $K$ algebraically closed, two $E$'s with the same $j$ are \unsure{isomorphic.} It is also easy to compute the group of automorphisms (but not, alas, the ring of endomorphisms). See Serre's ``Courbes Elliptiques - Formulaire'' for the details.
\endgroup

% archive sheet number: 4
\scanpage{6}\mspage{3.}
3. \uline{The} ``\uline{minimal}'' \uline{Weierstrass} \uline{equations} \uline{over} \uline{a} \uline{valuation} \uline{ring}. Let $v$ be a discrete valuation of $K$, with valuation ring $R$, prime ideal $\unsure{\mathfrak{p}|\wp}$, and residue field $k=R/\mathfrak{p}$. Let $E$ be an elliptic curve over $K$, with a $K$\unsure{.|-}rational point $O$.

\uline{Defn 3.1}: An equation for $E$ of the form (1.1) is minimal (with respect to $v$) if $v(a_{\unsure{i|1}})\geqslant0$ for all $i$ and if $v(\Delta)$ is minimal, subject to that condition.

\uline{Theorem 3.2}\unsure{.|:} \uline{A} \uline{minimal} \uline{equation} \inserted{for $E$} \uline{exists}, \uline{and} \uline{is} \uline{unique} \uline{up} \uline{to} \uline{a} \uline{change} \uline{of} \uline{coördinates} \uline{of} \uline{the} \uline{form} (2.1) \uline{with} $r,s,t\in R$ \uline{and} $u$ \uline{invertible} \uline{in} $R$.

Existence is obvious. Let $y^{\prime\,2}+a_1'x'y'+\ldots$ and $y^2+a_1xy+\ldots$ be two minimal equations for the same $E$. Since $\Delta\ne0$ and $v(\Delta')=v(\Delta)$ we conclude from\deleted{\illegible{1}} (2.3) that $v(u)=0$. Now from the transformation of $b_{\unsure{8|4}}$ and $b_6$ in (2.2) we see that $3r\in R$ and $4r\in R$, hence $r\in R$. \inserted{\unsure{Next|Now}} The transformation of $a_2$ shows that $s\in R$, and that of $a_6$ shows that $t\in R$.

\uline{Corollary 3.3:} \uline{The} \uline{differential} $\omega$ \uline{associated} \uline{with} \uline{a} \uline{minimal} \deleted{Weierstrass form} \uline{is} \uline{unique} \uline{up} \uline{to} \uline{a} \uline{unit} \uline{in} $R$.

\uline{Remarks}: 1) If $a_i\in R$ and $v(\Delta)<12$, then the equation (1.1) is automatically minimal. The converse is true if $j\in R$ and $p=\operatorname{char}k\ne2,3$. A complete algorithm for reducing to minimal form in all \deleted{\unsure{c}}\unsure{cases} is given below.

2) Let $\mathscr{M}$ be a collection of discrete valuations of $K$. For each $v\in\mathscr{M}$, let $y_v^2+a_{1,v}x_vy_v+\ldotp\,\ldotp\,\ldotp\,\ldotp$ be a minimal equation for $E$ relative to $v$, with corresponding differential $\omega_v$ and discriminant $\Delta_v$. Then the ``divisor'' $\displaystyle \unsure{\mathfrak{d}|\mathscr{D}}_E=\sum_{v\in\mathscr{M}}v(\Delta_v)v$ should be regarded as \uline{the} \uline{discriminant} \uline{of} $E$.

Let $F(x,y)=y^2+a_1xy+\ldots$ be an arbitrary Weierstrass equation for $E$ over $K$, and let $\omega_F$ and $\Delta_F$ be the corresponding differential and discriminant. Then the \uline{class} of the divisor $\displaystyle \unsure{\mathfrak{u}|\mathfrak{w}|u}_F=\sum_v v\left(\frac{\omega}{\omega_v}\right)v$ is independent of $F$, and we have $12\mathfrak{u}_F=\mathfrak{d}_E-(\Delta_F)$, so that $12\mathfrak{u}_F\sim\mathfrak{d}_E$.

3) In case $\mathscr{M}$ is the set of valuations associated with a principal ideal domain $D$ with field of fractions $K$, then it is easy to see that we can find one equation $F$ which simultaneously minimal for all $v$, so that $\mathfrak{u}_F=0$, and $\mathfrak{d}=(\Delta_F)$.

4) If $j\in R$, then $\Delta$ divides $c_4^3$ and $c_6^2$. From (1.7) we see that if $48\mathfrak{p}^4\mid c_4$ and $864\mathfrak{p}^6\mid c_6$, then the equation is not minimal. Hence \uline{if} $j\in R$ \uline{and} \uline{equation} \uline{minimal} \uline{we} \uline{have} $v(\Delta)<$ \deleted{\illegible{2}} $12+12v(2)+6v(3)$

% archive sheet number: 5
\scanpage{7}\mspage{4.}
\begingroup
\makeatletter\tagsleft@true\makeatother
4. \uline{The} \uline{canonical} \uline{filtration} \uline{on} \uline{the} \uline{group} of $v$-\uline{adic} \uline{points}. Let $F(x,y)\unsure{-|=}y^2+a_1xy+\unsure{\ldotp|\ldots}$ be a minimal equation for $E$ relative to a valuation $v$. Let $\widetilde F(x,y)$ be the reduction of $F$ $(\bmod\,\unsure{\mathfrak{p}|\wp})$ and let $\widetilde E$ denote the plane cubic $\widetilde F=0$ defined over the residue field $k$. By Theorem 3.2, $\widetilde E$ is uniquely determined by $E$ up to a projective transformation of the form (2.1) over $k$: Let $\widetilde E_0$ denote the non-singular part of $\widetilde E$. Then $\widetilde E_0$ is an algebraic group with origin $\widetilde O$. (If $\widetilde E$ is non singular, then $\widetilde E_0\unsure{-|=}\widetilde E$ is an elliptic curve. If $\widetilde E$ has a node
\parbox[c]{1.4em}{\setlength{\topsep}{0pt}\figure{\begin{tikzpicture}[x=0.95em,y=0.95em]
\draw[line width=.8pt,line cap=round] (0.97,.01) .. controls (.73,.19) and (.45,.66) .. (.24,.64) .. controls (.01,.65) and (-.02,.30) .. (.12,.16) .. controls (.29,.02) and (.49,.18) .. (.60,.34) .. controls (.76,.52) and (.85,.69) .. (.96,.87);
\end{tikzpicture}}}, then $\widetilde E_0\approx\mathbb{P}^1-(\text{two points})$ is a multiplicative group, and if $\widetilde E$ has a cusp
\parbox[c]{1.4em}{\setlength{\topsep}{0pt}\figure{\begin{tikzpicture}[x=0.95em,y=0.95em]
\draw[line width=.9pt,line cap=round] (.91,.66) .. controls (.81,.40) and (.59,.24) .. (.02,.22) .. controls (.61,.23) and (.87,.25) .. (1.01,.02);
\end{tikzpicture}}}, then $\widetilde E_0=\mathbb{P}^1-(\text{one point})$ is an additive group; here We have ignored questions of rationality, but if $k$ is perfect, so that the singularity of $\widetilde E$ is rational over $k$, then the analysis is the same, except that in case of a node, $\widetilde E_0$ is a multiplicative group ``twisted'' by the quadratic \unsure{extension} obtained by adjoining to $k$ the two tangents at the node)

Let $E(K)$ denote the group of points on $E$ rational over $K$, and let $\rho:E(K)\longrightarrow\widetilde E(k)$ denote the reduction map (defined naively in terms of the given projective coordinates --- by Theorem 3.2 it is independent of the coörds.) Let $E_0(K)=\rho^{-1}(\widetilde E_0(k))$ be the set of points whose reduction is non-singular.

\uline{Theorem 4.1}: $E_0(K)$ \uline{is} \uline{a} \uline{subgroup} \uline{of} \uline{finite} \uline{index} \uline{in} $E(K)$, \uline{and} $\rho_{\unsure{0}}:E_0(K)\to\widetilde E_0(k)$ \uline{is} \uline{a} \uline{homomorphism} \uline{of} \uline{groups}.

A straightforward but tedious proof can be given, using the addition formulae, for everything except the \unsure{`|``}finite index''. That finiteness depends on the minimality of the equation, and a proof of finiteness is implicit in the algorithm for reducing to minimal form given below.

We denote the kernel of $\rho_0$ by $E_1(K)$; it consists of the points $P=(x,y)$ in $E(K)$ such that $v(x)<0$ and $v(y)<0$. Clearly, from \unsure{(1.1)|(1 1)} since $v(a_{\unsure{i|1}})\geqslant0$, we have $v(x)<0\Longleftrightarrow v(y)<0$, in which case $v(x)=-2m$ and $v(y)\unsure{-|=}-3m$ for some $m$. For each $m\geqslant1$ we let $E_m(K)=\{(x,y)\in E(K)\mid v(x)\leq-2m\ \text{and}\ v(y)\leq\unsure{-3m}\}$ (understanding of course that $O\in E_m(K)$ for all $m$).

\uline{Theorem 4.2}: \uline{Let} $z=-\dfrac{x}{y}$. \uline{Then} $z$ \uline{is} \uline{a} \uline{uniformising} \uline{parameter} \uline{at} $O$. \uline{The} \uline{expansions}
\begingroup\scriptsize
\[
\begin{aligned}
x&=z^{\unsure{-|\cdot}2}-a_1z^{-1}-a_2-a_3z-(a_4+a_1a_3)z^2
 -\ldotp\,\ldotp\,\ldotp\,\ldotp\\
y&=-z^{-1}x=-z^{-3}+a_1z^{-2}+a_2z^{-1}+a_3+(a_4+a_1a_3)z+\ldots\\
\omega&=dt\bigl(1+a_1t+(a_1^2+a_2)t^2+(a_1^3+2a_1a_2+a_3)t^3+
 (a_1^4+3a_1^2a_2+6a_1a_3+a_2^2+2a_4)t^4+\unsure{\ldots}
\end{aligned}
\tag*{(4.3)}
\]
\endgroup
\uline{have} \uline{coefficients} \uline{in} $R$. \uline{So} \uline{also} \uline{does} \uline{the} \uline{formal}
% An earlier mathematical line runs underneath the two lines beginning here.
\deleted{$(\unsure{a_1^4+3a_1^2a_2}+6a_1\unsure{a_3}\unsure{+a_2^2+2a_4})\unsure{t^4+\ldots}$}
\uline{group} \uline{law} $\Phi(Z_1,Z_2)\unsure{=|:}Z_1+Z_2+\ldotp\,\ldotp\,\ldotp\,\ldotp$ \uline{defined} \uline{by} \uline{the} \uline{equation} $z(P+Q)=\Phi(z(P),z(Q))$. \uline{If} $R$ \uline{is} \uline{complete}, \uline{then} \uline{the} \uline{map} $z\longmapsto(x(z),y(z))$, \uline{for} $z\in\mathfrak{p}$, \uline{gives} \uline{an} \uline{isomorphism} \uline{of} $\mathfrak{p}_{\Phi}$ (\uline{the} \uline{prime} \uline{ideal} \uline{endowed} \uline{with} \uline{group} \uline{structure} \uline{via} $\Phi$) \uline{onto} $E_1(K)$, \uline{under} \uline{which} \uline{the} \uline{subgroups} $\mathfrak{p}^m_{\Phi}$ \uline{correspond} \uline{to} $E_m(K)$ \uline{for} \uline{all} $m\geqslant1$.
\endgroup

% archive sheet number: 6
\scanpage{8}\mspage{5}
\begingroup
\makeatletter\tagsleft@true\makeatother
The proof is straightforward. Let $z=\dfrac{x}{-y}$, $w=\dfrac{1}{-y}$, so that $x=\dfrac{z}{w}$ and $y=-\dfrac{1}{w}$. Then in terms of $w$ and $z$ the equation for $E$ is
\[
 w(1-a_1z-a_{\unsure{3|2}}w)\unsure{-|=}z^3+a_2z^2w+a_4zw^2+a_6w^3.
 \tag*{(4.4)}
\]
This shows that we have
\begingroup\scriptsize
\[
\begin{aligned}
w&=z^3+a_1z^4+(a_1^2+a_2)z^5+(a_1^3+2a_1a_2+a_3)z^6+
 (a_1^4+3a_1^2a_2+3a_1a_3+a_2^{\unsure{2|7}}+a_4)z^{\unsure{7|2}}+\ldots\\
 &\unsure{-|=}z^3+A_1z^4+A_2z^5+\ldots
\end{aligned}
\tag*{(4.5)}
\]
\endgroup
where $A_{\unsure{\nu|v}}$ is a polynomial of weight $\nu$ in the $a_i$ with positive integral coeffs\unsure{.} Hence the expansions of $y=-\dfrac{1}{w}$ and $x=-zy$ in terms of $z$ have coefficients in $\unsure{\mathbb{Z}|Z}[a_1,a_2,a_3,a_4,a_6]$. Now
\[
 \frac{\omega}{d\deleted{\unsure{t}}z}=
 \frac{dx/d\deleted{\unsure{t}}z}{2y+a_1x+a_3}=
 \frac{-2\deleted{\unsure{t}}z^{-3}+\ldotp}{-2\deleted{\unsure{t}}z^{-3}+\ldotp}=
 \frac{dy/d\deleted{\unsure{t}}z}{3x^2+2a_{\unsure{2}}x+a_{\unsure{4}}-a_1y}=
 \frac{3z^{-4}+\ldotp}{3z^{-4}+\ldots}
\]
has coefficients in $\mathbb{Z}[\frac{1}{2},a_1,\ldots,a_6]$ but also in $\mathbb{Z}[\frac{1}{3},a_1,\unsure{\ldotp},a_6]$, hence in $\mathbb{Z}[a_1,\unsure{\ldotp},a_6]$. As for the group law, if $z_1$ and $z_2\in\unsure{\mathfrak{p}|\wp}$, then the line joining the points $(z_1,w_1)$, $(z_{\unsure{1|2}},w_2)$ in the $(z,w)$-plane has slope $\unsure{\in}\mathfrak{p}^2$, because
\[
 \frac{w_2-w_1}{z_2-z_1}=
 \frac{z_2^3-z_1^3}{z_2-z_{\unsure{1}}}
 +A_1\frac{z_2^4-z_1^4}{z_2-z_1}+
\]
\margin{\unsure{with} $A_i$ as above in (\unsure{4}.5)}
Call this slope $\lambda=\lambda(z_1,z_2)\unsure{-|=}z_2^{\unsure{2|7}}+z_1z_2+z_1^2+A_1(z_2^3\unsure{+|-}z_2^{\unsure{1|2}}z_1+z_{\unsure{1|2}}z_1^{\unsure{2|7}}+z_{\unsure{1}}^3)+\ldots$. Put $\nu=\nu(z_1,z_2)=w_i-\lambda z_i$ $(i=1,2)$. Substituting $w=\lambda(z_1,z_2)z+\nu(z_1,z_2)z$ in (4.4) we find a cubic in $z$ with roots $w_1$ and $w_2$. Looking at the sum of the roots, one sees that the third root $\deleted{\unsure{t}}z_3$ is expressed as a power series in $z_1,z_2$ with coeffs in $R$. In fact,
\[
 z_1+z_2+z_3=
 \frac{a_1\lambda+a_3\lambda^2-a_2\nu-2a_4\lambda\nu-3a_6\lambda^{\unsure{2|7}}\nu}
 {1+a_2\lambda+a_4\lambda^2+a_6\lambda^{\unsure{3}}}
 \tag*{(4.6)}
\]
\margin{(sauf erreur)}
Thus we have the \unsure{`|``}canonical filtration''
\[
 E(K)\supset E_0(K)\supset E_1(K)\supset E_2(K)\supset\ldots
 \qquad\supset\unsure{(O)|(0)}=\bigcap_{m\unsure{=|-}1}^{\infty}E_m(K)
 \tag*{(4.7)}
\]
with quotients
\begingroup\small
\[
 E(K)\big/E_0(K)\quad\text{finite},\qquad
 E_0\big/E_1\mathrel{\deleted{\unsure{\subset}}\hookrightarrow}\widetilde E_0(k),
 \quad\text{and}\quad E_m\big/E_{m+1}\hookrightarrow k^+
 \quad\text{for }m\geqslant1.
\]
\endgroup
Of course the inclusions $\hookrightarrow$ are bijections $\overset{\sim}{\to}$ if $R$ is complete.
\endgroup

% archive sheet number: 7
